\scanpage{11}
\uline{Théorème} \inserted{Soit $k$ un corps \unsure{du|de} \unsure{Thm.}} Soit $K$ un groupe algébrique sur $k$, et $K'$ un noyau sur $k$, localement isomorphe à $K$. Alors $H^2(k,K')=H^2(k,K')'=\deleted{\unsure{\{e\}}}$ uniponctuel.

En effet, $K'$ provient d'un élément de $H^1(k,\underline{\operatorname{Aut}\operatorname{Ext}}(K))$, qui par le raisonnement précédent [\uline{légèrement généralisé}, comme $\underline{\operatorname{Aut}}(K)$ n'est pas néc. un groupe algébrique\unsure{\ldots}] provient d'un élément de $H^1(k,\underline{\operatorname{Aut}}(K))$, ce qui prouve que le noyau $K'$ est représentable par un groupe $K_1$, [loc. isom. à $K$], \deleted{d'où déjà} \inserted{i.e.} \deleted{$\unsure{H^2(k,}$} $H^2(k,K')'\ne\varnothing$. Comme alors $H^2(k,K')\ne\varnothing$ est un fibré principal homogène sous $H^2(k,\unsure{\mathfrak{z}|Z})$, $\mathfrak{z}$ le centre de $\unsure{K_1|K'}$, ou encore de $K$, or on sait que $H^2(k,\mathfrak{z})=0$. Donc $H^2(k,K')$ réduit à un \unsure{pt}, donc $H^2(k,K')'$ aussi [étant non vide et \unsure{dedans contenu}].
% archive sheet number: 10

